\section{Vector Database 与相似度搜索}

Vector Database 是任何严肃记忆系统的核心。与 SQL 的精确匹配不同，它在高维空间中找到一个向量的最近邻居，实现\textbf{语义搜索}——即使没有共享词汇也能找到概念相关的记忆。

\subsection{相似度搜索原理}

每条记忆被转换为一个向量（以 OpenAI Embedding 模型为例，是一个包含 1,536 个浮点数的数组）。概念相似的文本会产生相似的向量。查询时，将查询文本也嵌入为向量，然后通过 Cosine Similarity 找到最近的向量。

$$
\text{cosine\_similarity}(\mathbf{a}, \mathbf{b}) = \frac{\mathbf{a} \cdot \mathbf{b}}{\|\mathbf{a}\| \cdot \|\mathbf{b}\|}
$$

其中：
\begin{itemize}[leftmargin=2em]
    \item $\mathbf{a}, \mathbf{b}$：两个 Embedding 向量
    \item $\mathbf{a} \cdot \mathbf{b}$：向量点积
    \item $\|\mathbf{a}\|, \|\mathbf{b}\|$：向量的 L2 范数
\end{itemize}

\begin{knowledgebox}{Vector Database 选型建议}
\begin{itemize}[leftmargin=1.5em]
    \item \textbf{ChromaDB}：本地开发首选，轻量易上手
    \item \textbf{pgvector}：如果你已在用 PostgreSQL，零额外基础设施
    \item \textbf{Pinecone / Qdrant}：需要大规模部署时使用
\end{itemize}
\end{knowledgebox}

\subsection{本章小结}

Vector Database 通过将文本转为高维向量并计算 Cosine Similarity 来实现语义搜索。选择工具时应根据开发阶段和规模需求决定。

